\begin{center}
    \textbf{\large Вариант 2}
\end{center}

\newpage
\null
\newpage

\samerowans{\begin{problem}{1}
Вычислите: $7567 - 7245 : (416 - 15 \cdot 17) + 78$
\end{problem}}

\vspace{5pt}

\samerowans{\begin{problem}{2}
Решите уравнение: $(24 \cdot 4 - 252 : b) \cdot 4 = 48$
\end{problem}}

\vspace{5pt}

\samerowans{\begin{problem}{3}
Вычислите: $8\text{ т } 45\text{ кг } - 54\text{ ц } + 136\text{ кг}$. Ответ запишите в т, ц, кг.
\end{problem}}

\samerowans{\begin{problem}{4}
Сложили пять одинаковых чисел, к результату прибавили 3, затем полученную сумму разделили на 9 и получили 17. Какие числа складывали?
\end{problem}}

\vspace{5pt}

\samerowans{\begin{problem}{5}
Юля два часа делала домашнее задание. Две трети часа она делала математику, полчаса русский. Все оставшееся время она учила английский. Сколько минут Юля учила английский?
\end{problem}}

\vspace{5pt}

\begin{problem}{6} Из двух пунктов А и В, расстояние между которыми 40 км, одновременно навстречу друг другу выехали два велосипедиста. Скорость велосипедиста, выехавшего из п А, на $5\text{ км/ч}$ больше скорости второго. За два часа им удалось встретиться и разъехаться на расстояние 30 км. На каком расстоянии от п А будет второй велосипедист через 3 часа 20 минут после своего выезда из В?
\end{problem}

\vspace{5pt}

\noindent\grid{36}{19}

\vspace{5pt}

\samerowans{\begin{problem}{7}
Первый станок может сделать 320 деталей за 8 часов, второй справится с этой задачей в 2 раза быстрее. Сколько часов совместной работы потребуется двум станкам, чтобы сделать 720 деталей?
\end{problem}}

\newpage

\begin{problem}{8} Из прямоугольника с периметром 20 дм вырезали квадрат со стороной 6 см. Найди площадь оставшейся фигуры, если длина прямоугольника 600 мм.
\end{problem}

\smallskip

\noindent\grid{36}{15}

\vspace{5pt}

\samerowans{\begin{problem}{9}
Таня заплатила за покупку 1 м 50 см шелковой ленты 750 рублей. Надя купила ленты на 40 см больше, чем Таня. Сколько рублей заплатила Надя?
\end{problem}}

\vspace{5pt}

\samerowans{\begin{problem}{10}
Коля и Петя сыграли партию в шахматы, которая закончилась победой одного из них. «Петя победил», — сказал Коля. «Коля проиграл», — сказал Петя. Кто из них выиграл, если известно, что хотя бы кто-то из них сказал неправду?
\end{problem}}

\vspace{5pt}

\begin{problem}{11} Через четыре года мальчик будет вдвое старше, чем он был четыре года назад. А девочка будет через три года втрое старше, чем семь лет назад. Кто старше: мальчик или девочка?
\end{problem}

\smallskip

\noindent\grid{36}{19}

\newpage

